\section{Сложение}

\begin{definition}{2.2.1}{сложение натуральных чисел}
Пусть $m$ — натуральное число. Прибавление нуля к $m$ определяем как $0+m := m$.
Предположим по индукции, что прибавление $n$ к $m$ уже определено. Тогда прибавление
$n{+}{+}$ к $m$ определяем как $(n{+}{+})+m := (n+m){+}{+}$.
\leansource{Nat.add}{(n m : Nat) : Nat}
\end{definition}

\begin{remark}
У нас есть только одна исходная операция — прибавление единицы — и несколько аксиом.
$3+5$ — это то же самое, что $((5{+}{+}){+}{+}){+}{+}$: прибавить $3$ к $5$ значит трижды
взять следующее число. Формально это рекурсивное определение из
\taoref{prop:2.1.16}{утверждения 2.1.16}, применённое к $c=m$ и $f_n(a) = a{+}{+}$. Оно
асимметрично по $n$ и $m$: индукция идёт по первому слагаемому, и оба равенства
$0+m=m$, $(n{+}{+})+m=(n+m){+}{+}$ верны по построению.
\end{remark}

\begin{lemma}{2.2.2}{}
Для любого натурального $n$ выполняется $n+0=n$.
\leansource{Nat.add_zero}{(n : Nat) : n + 0 = n}
\end{lemma}

\begin{proof}
\prooflabel{База} $0+0=0$, поскольку $0+m=m$ для любого $m$.\\
\prooflabel{Переход} Пусть $n+0=n$. По определению сложения $(n{+}{+})+0=(n+0){+}{+}$,
что по предположению индукции равно $n{+}{+}$.
\end{proof}

\begin{remark}
Это не следует напрямую из $0+m=m$: там ноль стоит слева, здесь справа, и без
коммутативности это два разных утверждения. Коммутативность же сама опирается на эту
лемму.
\end{remark}

\begin{lemma}{2.2.3}{}
Для любых натуральных $n$ и $m$ выполняется $n+(m{+}{+})=(n+m){+}{+}$.
\leansource{Nat.add_succ}{(n m : Nat) : n + (m++) = (n + m)++}
\end{lemma}

\begin{proof}
Индукция по $n$ при фиксированном $m$.\\
\prooflabel{База} $0+(m{+}{+})=m{+}{+}=(0+m){+}{+}$.\\
\prooflabel{Переход} Пусть $n+(m{+}{+})=(n+m){+}{+}$. Тогда по определению сложения
\[
  (n{+}{+})+(m{+}{+}) = (n+(m{+}{+})){+}{+} = ((n+m){+}{+}){+}{+},
\]
а по тому же определению $(n{+}{+})+m=(n+m){+}{+}$, откуда правая часть равна
$((n{+}{+})+m){+}{+}$ — то же самое.
\end{proof}

\begin{remark}
Отсюда и из \taoref{lem:2.2.2}{леммы 2.2.2} следует $n{+}{+}=n+1$: для любого $n$ «прибавить
единицу» и «взять следующее» — одно и то же.
\leansource{Nat.succ_eq_add_one}{(n : Nat) : n++ = n + 1}
\end{remark}

\begin{proposition}{2.2.4}{сложение коммутативно}
Для любых натуральных $n$ и $m$ выполняется $n+m=m+n$.
\leansource{Nat.add_comm}{(n m : Nat) : n + m = m + n}
\end{proposition}

\begin{proof}
Индукция по $n$ при фиксированном $m$.\\
\prooflabel{База} $0+m=m$ по определению, и $m+0=m$ по \taoref{lem:2.2.2}{лемме 2.2.2} —
обе части равны $m$.\\
\prooflabel{Переход} Пусть $n+m=m+n$. По определению сложения $(n{+}{+})+m=(n+m){+}{+}$.
С другой стороны, $m+(n{+}{+})=(m+n){+}{+}$ по \taoref{lem:2.2.3}{лемме 2.2.3}, что равно
$(n+m){+}{+}$ по предположению индукции.
\end{proof}

\begin{proposition}{2.2.5}{сложение ассоциативно}
Для любых натуральных $a,b,c$ выполняется $(a+b)+c=a+(b+c)$.
\leansource{Nat.add_assoc}{(a b c : Nat) : (a + b) + c = a + (b + c)}
\end{proposition}

\begin{proof}
Индукция по $a$ при фиксированных $b,c$.\\
\prooflabel{База} $(0+b)+c=b+c$ и $0+(b+c)=b+c$ по определению — обе части равны $b+c$.\\
\prooflabel{Переход} Пусть $(a+b)+c=a+(b+c)$. По определению сложения
$(a{+}{+})+b=(a+b){+}{+}$, откуда $((a{+}{+})+b)+c = ((a+b){+}{+})+c = ((a+b)+c){+}{+}$ —
снова по определению, применённому к первому слагаемому суммы $((a+b){+}{+})+c$. По
предположению индукции это равно $(a+(b+c)){+}{+}$, а это и есть $(a{+}{+})+(b+c)$ по
определению.
\end{proof}

\begin{proposition}{2.2.6}{закон сокращения}
Пусть $a,b,c$ — натуральные числа и $a+b=a+c$. Тогда $b=c$.
\leansource{Nat.add_left_cancel}{(a b c : Nat) (habc : a + b = a + c) : b = c}
\end{proposition}

\begin{proof}
Индукция по $a$ при фиксированных $b,c$.\\
\prooflabel{База} $0+b=0+c$ означает $b=c$ напрямую, по определению сложения.\\
\prooflabel{Переход} Пусть из $a+b=a+c$ следует $b=c$. Дано $(a{+}{+})+b=(a{+}{+})+c$, то есть
$(a+b){+}{+}=(a+c){+}{+}$. По \taoref{ax:2.4}{аксиоме 2.4} отсюда $a+b=a+c$, а по
предположению индукции — $b=c$.
\end{proof}

\begin{remark}
Вычитания и отрицательных чисел у нас пока нет, но этот закон уже даёт своего рода
«виртуальное вычитание»: он позволяет сокращать одинаковое слагаемое в обеих частях
равенства, не выражая явно $b=(a+c)-a$.
\end{remark}

\begin{definition}{2.2.7}{положительные натуральные числа}
Натуральное число $n$ называется \emph{положительным}, если $n \neq 0$.
\leansource{Nat.IsPos}{(n : Nat) : Prop := n ≠ 0}
\end{definition}

\begin{proposition}{2.2.8}{}
Если $a$ положительно, а $b$ — произвольное натуральное число, то $a+b$ положительно.
\leansource{Nat.add_pos_left}{(b : Nat) (ha : a.IsPos) : (a + b).IsPos}
\end{proposition}

\begin{proof}
Индукция по $b$ при фиксированном положительном $a$.\\
\prooflabel{База} $a+0=a$ по \taoref{lem:2.2.2}{лемме 2.2.2}, а $a$ положительно по
условию.\\
\prooflabel{Переход} Пусть $a+b$ положительно. Тогда $a+(b{+}{+})=(a+b){+}{+}$, а это не
может быть нулём по \taoref{ax:2.3}{аксиоме 2.3} — значит, тоже положительно.
\end{proof}

\begin{remark}
Вместе с \taoref{prop:2.2.4}{коммутативностью} отсюда сразу следует и $b+a$ положительно.
\end{remark}

\begin{corollary}{2.2.9}{}
Если $a+b=0$, то $a=0$ и $b=0$.
\leansource{Nat.add_eq_zero}{(a b : Nat) (hab : a + b = 0) : a = 0 ∧ b = 0}
\end{corollary}

\begin{proof}
Предположим противное: $a \neq 0$ или $b \neq 0$.\\
Если $a \neq 0$, то $a$ положительно, а тогда $a+b$ положительно.\by{\taoref{prop:2.2.8}{утверждение 2.2.8}}\\
Аналогично, если $b \neq 0$, то $b$ положительно, и снова $a+b$ положительно.\by{\taoref{prop:2.2.8}{утверждение 2.2.8}}\\
Оба случая противоречат условию $a+b=0$.
\end{proof}

\begin{lemma}{2.2.10}{единственный предшественник}
Пусть $a$ — положительное число. Тогда существует единственное натуральное $b$ такое, что
$b{+}{+}=a$.
\leansource{Nat.uniq_succ_eq}{(a : Nat) (ha : a.IsPos) : ∃! b, b++ = a}
\end{lemma}

\begin{proof}
Существование — индукцией по $a$: случай $a=0$ невозможен, так как $a$ положительно. Если
$a=n{+}{+}$, то $b:=n$ подходит.\\
Единственность: если $y_1{+}{+}=a=y_2{+}{+}$, то $y_1=y_2$ по \taoref{ax:2.4}{аксиоме 2.4}.
\end{proof}

\begin{definition}{2.2.11}{порядок на натуральных числах}
Пусть $n,m$ — натуральные числа. Говорим, что $n$ \emph{больше или равно} $m$, и пишем
$n \geq m$ или $m \leq n$, если $n=m+a$ для некоторого натурального $a$. Говорим, что $n$
\emph{строго больше} $m$, и пишем $n>m$ или $m<n$, если $n \geq m$ и $n \neq m$.
\leansource{Nat.instLE}{: LE Nat}
\leansource{Nat.instLT}{: LT Nat}
\end{definition}

\begin{proposition}{2.2.12}{базовые свойства порядка}
Пусть $a,b,c$ — натуральные числа. Тогда:
\begin{itemize}
  \item[(a)] (рефлексивность) $a \geq a$;
  \item[(b)] (транзитивность) из $a \geq b$ и $b \geq c$ следует $a \geq c$;
  \item[(c)] (антисимметричность) из $a \geq b$ и $b \geq a$ следует $a=b$;
  \item[(d)] (сложение сохраняет порядок) $a \geq b$ равносильно $a+c \geq b+c$;
  \item[(e)] $a<b$ равносильно $a{+}{+} \leq b$;
  \item[(f)] $a<b$ равносильно тому, что $b=a+d$ для некоторого положительного $d$.
\end{itemize}
\leansource{Nat.ge_refl}{(a : Nat) : a ≥ a}
\leansource{Nat.ge_trans}{(hab : a ≥ b) (hbc : b ≥ c) : a ≥ c}
\leansource{Nat.ge_antisymm}{(hab : a ≥ b) (hba : b ≥ a) : a = b}
\leansource{Nat.add_ge_add_right}{(a b c : Nat) : a ≥ b ↔ a + c ≥ b + c}
\leansource{Nat.lt_iff_succ_le}{(a b : Nat) : a < b ↔ a++ ≤ b}
\leansource{Nat.lt_iff_add_pos}{(a b : Nat) : a < b ↔ ∃ d : Nat, d.IsPos ∧ b = a + d}
\end{proposition}

\begin{remark}
Как и в книге, доказательства этих шести пунктов оставлены здесь в стороне — все они
прямые следствия \taoref{def:2.2.11}{определения 2.2.11} и уже доказанных свойств сложения
(в Lean все доказаны явно, см. ссылки выше).
\end{remark}

\begin{proposition}{2.2.13}{трихотомия порядка}
Пусть $a,b$ — натуральные числа. Тогда выполняется ровно одно из трёх: $a<b$, $a=b$
или $a>b$.
\leansource{Nat.trichotomous}{(a b : Nat) : (a < b) ∨ (a = b) ∨ (a > b)}
\end{proposition}

\begin{proof}
То, что верно не больше одного из трёх, — из \taoref{prop:2.2.12}{утверждения 2.2.12(c)}:
если бы, скажем, $a<b$ и $a>b$ одновременно, антисимметричность дала бы $a=b$, что
противоречит $a \neq b$ из определения $<$.\\
То, что верно хотя бы одно из трёх, — индукцией по $a$ при фиксированном $b$: база $a=0$
разбирается через \taoref{prop:2.2.12}{утверждение 2.2.12(f)} ($0 \leq b$ всегда), а на
шаге индукции трихотомия для $a$ подсказывает, какой из трёх случаев даёт трихотомию для
$a{+}{+}$.
\end{proof}

\begin{proposition}{2.2.14}{усиленный принцип индукции}
Пусть $m_0$ — натуральное число, и пусть $P(m)$ — свойство, применимое к натуральным
числам $m \geq m_0$. Предположим, что для каждого $m \geq m_0$ из истинности $P(m')$ при
всех $m_0 \leq m' < m$ следует истинность $P(m)$. Тогда $P(m)$ истинно для всех
$m \geq m_0$.
\leansource{Nat.strong_induction}{(hind : ∀ m, m ≥ m₀ → (∀ m', m₀ ≤ m' ∧ m' < m → P m') → P m) : ∀ m, m ≥ m₀ → P m}
\end{proposition}

\begin{remark}
Книга, как и здесь, оставляет доказательство в стороне — оно сводится к обычной
\taoref{ax:2.5}{аксиоме 2.5}, применённой к вспомогательному свойству «$P$ верно для всех
$m'$ от $m_0$ до $m$». На практике чаще всего берут $m_0=0$ или $m_0=1$. Родственный
вариант — \emph{обратная индукция} (\code{Nat.backwards\_induction}): если $P(n)$ истинно
и $P(m{+}{+})$ влечёт $P(m)$, то $P(m)$ истинно для всех $m \leq n$.
\end{remark}
